\documentclass{article}
\usepackage{tikz}
\usetikzlibrary{angles, quotes}
\usepackage{amssymb}
\usepackage{pgfplots}
\pgfplotsset{width=10cm,compat=1.9}
\usepgfplotslibrary{statistics}
\usepackage[left=2.54cm,right=2.54cm,top=2.54cm,bottom=2.54cm]{geometry}
\usepackage{tasks} %Paquete necesario para la producción de items horizontales para algunos ejercicios de tarea empleando el comando /begin{multicols}}{#}.../end{multicols} para ayudarnos a reducir las páginas
\usepackage{pdflscape} %comando que permite cambiar a orientación horizontal las hojas%
\usepackage{soul} %Paquete para permitir la compilación de acentos usando el comando {\hl{}}}\\
\sethlcolor{yellow(munsell)} %comando para definir el color para el subrayado usando el comando \hl
\usepackage[utf8]{inputenc}
\usepackage[spanish]{babel}
\usepackage{icomma}
\usepackage{ragged2e}
\usepackage{siunitx}
\usepackage{url}
\usepackage[colorlinks=true, urlcolor=blue,  linkcolor=black, citecolor=green]{hyperref}
\usepackage{pdfpages}
\usepackage{blindtext} %Paquete que genera texto ficticio (dummy text) con el comando: \blindtext
\usepackage{cancel} %in the preamble gives you four different modes of striking through: \cancel{text to cancel} draws a diagonal line (slash) through its argument, \bcancel{text to cancel} uses the negative slope (a backslash), \xcancel{text to cancel} draws an X (actually \cancel plus \bcancel), \cancelto{〈value〉}{〈expression〉} draws a diagonal arrow through the 〈expression〉pointing to the 〈value〉 (math-mode only)
\usepackage{csquotes}
\usepackage{afterpage}
\usepackage{parskip} 
\usepackage{float}
\usepackage{enumitem}
\usepackage{multicol}%Paquete que permite la creación aislada de columnas de texto. De esta forma se reduce la cantidad de páginas en nuestro documento
\usepackage{lipsum} %Paquete que genera texto de relleno al igual que \usepackage{blindtext}, se usa con el comando \lipsum[#-#] los #(hashtags) delimitan el número de páginas que deseamos ocupar del paquete lipsum para usarlos como relleno. 

% Definir el comando personalizado para la nota
\newcommand{\mynote}[1]{%
\begin{tikzpicture}[baseline=-0.75ex]
\node[draw, circle, fill=Apple Green!60, inner sep=2pt] (note) {\textbf{Nota:}};
\end{tikzpicture}%
\ \textit{#1}%
}

\newenvironment{Figura}
{\par\medskip\noindent\minipage{\linewidth}}
{\endminipage\par\medskip}
\usepackage{caption}
\usepackage[
backend=biber,
sorting=none,
url=true
]{biblatex} %bibliografía
\addbibresource{biblio.bib}
% Establecer el espacio entre las entradas de la bibliografía
\setlength{\bibitemsep}{1\baselineskip} % Puedes ajustar el valor según tus preferencias
\usepackage{amsmath,amsthm,amssymb,amsfonts}
\usepackage{pifont} %Permite la compilación del comando \xmark: es el símbolo de la tachita
\usepackage{mathtools} %permite la compilación de símbolos para matrices
\usepackage{empheq} %Paquete que se relaciona con \usepackage[most]{tcolorbox} para la creación de cuadros/cajas de colores para delimitar los resultados matemáticos o líneas de texto usando la línea de comando: \begin{empheq}[box={\mymath[colback=orange(webcolor)!70,drop lifted shadow]}]{equation*} \end{empheq}
\usepackage[most]{tcolorbox} %Paquete que se relaciona con \usepackage{empheq} para la creación de cuadros/cajas de colores para delimitar los resultados matemáticos o líneas de texto
\renewcommand{\qedsymbol}{$\blacksquare$}
\usetikzlibrary{positioning,decorations.pathreplacing} %Paquete que permite la compilación de llaves para cuadros sinópticos (brace diagram) 
\usepackage{schemata} %The schemata package is designed just for these brace diagrams
\usepackage{booktabs}

\newcommand\AB[2]{\schema{\schemabox{#1}}{\schemabox{#2}}} %comando o paquete necesario para crear cuadros sinópticos

\newtcbox{\mymath}[1][]{%
nobeforeafter, math upper, tcbox raise base,
enhanced, colframe=blue!30!black,
colback=blue!30, boxrule=1pt,
#1} %Comando importante para encasillar los resultados matemáticos en cajas de diferentes colores, se relaciona con los paquetes \usepackage{empheq} y \usepackage[most]{tcolorbox} 

\newenvironment{sysmatrix}[1]
{\left(\begin{array}{@{}#1@{}}}
{\end{array}\right)}
\newcommand{\ro}[1]{%
\xrightarrow{\mathmakebox[\rowidth]{#1}}%
}
\newlength{\rowidth}% row operation width
\AtBeginDocument{\setlength{\rowidth}{3em}}  %Comando importante para laproducción de lineas de operaciones entre matrices, método de Gauss o Gauss Jordan se relaciona con \newenvironment{sysmatrix}

%formato para cambiar el horario a español
\usepackage[spanish]{datetime2}
\DTMsetdatestyle{spanish}

\renewcommand{\today}{\DTMdisplaydate{\the\year}{\the\month}{\the\day }{-1}}
%%%%%%%%%%%%%%%%%%%%%%%%%%%%%%%%%%%%%%%%%

\usepackage{datetime}
\newcommand{\mycurrenttime}{\xxivtime}

%%%%%%%%%%%%%%%%%%%%%%%%%%%%%%%%%%%%%%%%%%%%%%%%%%%%%%%%%%%%%%%%%%%%%%%%%%%%%%%%%%%%%%%%%%%%%%%%%%%%%%%%%%%%%%%%%%%%%%%%%%%%%%%%%%%%%%%%%%%%%%%%%


%%%%%%%%%%%%%Caja de color para el título%%%%%%%%%%%%%%%%%%%%%%%%%%%%%%%%%%%%%%%%%%%%%%%%%%%%%%

\definecolor{myframecolor}{RGB}{1, 45, 75} % Prussian Blue
\definecolor{myboxcolor}{RGB}{0, 107, 92} % Apple Green

\newtcolorbox{mybox}{
enhanced,
colback=myboxcolor!7, % Color de fondo del cuadro
colframe=myframecolor, % Color del marco del cuadro
arc=0pt,
boxrule=1pt,
borderline west={2mm}{-10mm}{myframecolor}, % Borde en el lado izquierdo
sharp corners=southwest,
width=\linewidth,
before=\par\vspace{\bigskipamount}, % Espacio antes del cuadro
after=\par\vspace{\bigskipamount} % Espacio después del cuadro
}


\newcommand{\euler}{e} %Comando para producir letra e de euler
\newenvironment{remark}{\par\vfill\footnotesize % Vertical white space above the remark and smaller font size
\begin{list}{}{
	\leftmargin=80pt % Indentation on the left
	\rightmargin=60pt}\item\ignorespaces % Indentation on the right
\makebox[-2.5pt]{\begin{tikzpicture}[overlay]
		\node[draw=Horizon!60,line width=2.5pt,circle,fill=Horizon!25,font=\sffamily\bfseries,inner sep=4pt,outer sep=0pt] at (-19pt,5pt){\textcolor{Horizon}{Nota}};\end{tikzpicture}} % Blue Nota in a circle
\advance\baselineskip -1pt}{\end{list}\vskip5pt} % Tighter line spacing and white space after remark
\usepackage{graphicx}

\usepackage{titling}

\input{colores.tex}

%%Producir signos de cita textual``Asignatura''

\newcommand{\xmark}{\ding{55}} %%%comando que genera la tachita

\newenvironment{MyColorPar}[1]{%
\leavevmode\color{#1}\ignorespaces%
}{%
}%

%%%%%%%%%%%%%%%%%%% Título de la tarea, Nombre de alumno

\title{Clase \# 31 - Mecánica Teórica}
\author{\emph{JA} $\boldsymbol{\mid}$}
\date{\today}

%%%%%%%%%%%%%%%%%%% Datos de la Materia

\usepackage{fancyhdr}
\fancypagestyle{plain}{%  the preset of fancyhdr 
\fancyhf{} % clear all header and footer fields
\fancyfoot[R]{
\includegraphics[width=2cm]{LogoFCFMBUAP (1).png}}
\fancyfoot[L]{{\bfseries{\thedate{}}} a las {\bfseries{\mycurrenttime{}}} horas{} (GMT-6, H. Puebla de Zaragoza, Pue)}
\fancyhead[L]{Mecánica teórica}
\fancyhead[R]{\theauthor}
}
\makeatletter
\def\@maketitle{%
\newpage
\null
\vskip 1em%
\begin{center}%
\let \footnote \thanks
{\LARGE \@title \par}%
\vskip 1em%
%{\large \@date}%
\end{center}%
\par
\vskip 1em}
\makeatother

\usepackage{lipsum}  
\usepackage{cmbright}


%%%%%%%%%%%%%%%%%%%%%%%%%%%%%%%%%%%%%%%%%%%%%%%%%%%%%%%%%%%%%%%%%%%%%%%%%%%%%%%%%%%%%%%%%%%%%%%%%%%%%%%%%%%%%%%%%%%%%%%%%%%%%%%%%%%%%%%%%%%%%%%%%%%%%%%%%%%%%%%%%%%%%%%%%%%%%%%%%%%%%%%%%%%%%%%%%%%%%%

\begin{document}



\maketitle


%%%%%%%%%%%%%%%%%%% Datos del profesor, campus y aula

\begin{mybox}
\noindent\begin{tabular}{@{}ll}
	{\bfseries{Profesor}} & Dr. Gilberto Silva Ortigoza\\
	\textcolor{prussianblue}{\bfseries{Campus}} / \textcolor{trueblue}{\bfseries{Facultad}}  & \textcolor{prussianblue}{\bfseries{C.U. BUAP}} /  \textcolor{trueblue}{\bfseries{FCFM}} \\
	\textcolor{prussianblue}{\bfseries{Edificio}} /  \textcolor{trueblue}{\bfseries{EMA4}}    & 	\textcolor{prussianblue}{\bfseries{FM9}} /  \textcolor{trueblue}{\bfseries{302}}   
\end{tabular} 
\end{mybox} \vspace{1cm}

\tableofcontents

\vspace{2cm}

\section{Próximo examen: problema de Kepler (dos cuerpos)}

\subsection*{Qué dominar}
\begin{itemize}
  \item Identificar el \textbf{potencial kepleriano}: \(V(r)=-\dfrac{\kappa}{r}\), con \(\kappa=G\,m_1 m_2\).
  \item Escribir el \textbf{lagrangiano} y separar el \textbf{movimiento del centro de masa} y el \textbf{movimiento relativo}.
  \item Derivar las \textbf{ecuaciones de movimiento} para una \textbf{fuerza central} y sus consecuencias: conservación de \(\vec{L}\), movimiento planar y forma de la órbita.
\end{itemize}

\subsection*{Lagrangiano y separación CM/relativo}
Sea
\[
M=m_1+m_2,\qquad \mu=\frac{m_1 m_2}{m_1+m_2},\qquad \kappa=G\,m_1 m_2.
\]
Con \(\vec{R}=\dfrac{m_1\vec{r}_1+m_2\vec{r}_2}{M}\) (centro de masa) y \(\vec{r}=\vec{r}_2-\vec{r}_1\) (coordenada relativa), el lagrangiano se separa como
\[
L=\underbrace{\frac{1}{2}M\,\dot{\vec{R}}^{\,2}}_{\text{coordenadas mayúsculas }(CM)}
\;+\;
\underbrace{\Big(\frac{1}{2}\mu\,\dot{\vec{r}}^{\,2}+\frac{\kappa}{r}\Big)}_{\text{coordenadas minúsculas }(\text{relativo})}.
\]
Es decir,
\[
L=L_{\vec{R}}+L_{\vec{r}},\qquad
L_{\vec{R}}=\frac{1}{2}M\,\dot{\vec{R}}^{\,2},\qquad
L_{\vec{r}}=\frac{1}{2}\mu\,\dot{\vec{r}}^{\,2}+\frac{\kappa}{r}\quad (V=-\kappa/r).
\]

\begin{figure}[H]
    \centering
	    \includegraphics[width=0.8\textwidth]{1.png} % ajusta el ancho según convenga
    \caption{Descripción clara y concisa de la figura.}
    \label{fig:mi_figura}
\end{figure}

\subsection*{Centro de masa}
Del principio de inercia del CM:
\[
\frac{d}{dt}\big(M\dot{\vec{R}}\big)=\vec{0}
\;\;\Rightarrow\;\;
\vec{R}(t)=\vec{R}_{0}+\vec{V}_{0}\,t
\quad\text{(MRU).}
\]

\subsection*{Movimiento relativo y fuerza central}
Para el movimiento relativo,
\[
\mu\,\ddot{\vec{r}}=-\frac{\partial V}{\partial r}\,\hat{\mathbf r}=-\frac{\kappa}{r^2}\,\hat{\mathbf r}
\;\;\Leftrightarrow\;\;
\mu\,\ddot{\vec{r}}=-\kappa\,\frac{\vec{r}}{r^3}.
\]
Al ser una \textbf{fuerza central} \(\vec{F}=f(r)\,\hat{\mathbf r}\), el \textbf{momento angular orbital} se conserva:
\[
\vec{L}=\mu\,\vec{r}\times\dot{\vec{r}}=\text{cte},\qquad
\dot{\vec{L}}=\vec{r}\times\vec{F}=\vec{0}.
\]
Esto fija un \textbf{plano de movimiento} perpendicular a \(\vec{L}\). La “sabor esférico” del problema proviene de la simetría radial del potencial (depende solo de \(r\)), pero la dinámica efectiva vive en un plano con coordenadas polares \((r,\phi)\).

\begin{figure}[H]
    \centering
	    \includegraphics[width=0.6\textwidth]{2.png} % ajusta el ancho según convenga
    \caption{Descripción clara y concisa de la figura.}
    \label{fig:mi_figura2}
\end{figure}

\subsection*{Geometría angular y ortogonalidad}
Como \(\vec{L}\) es constante, podemos tomar un marco inercial con \(\vec{L}=\ell\,\hat{\mathbf z}\). Se cumple
\[
\vec{L}\perp\vec{r},\qquad \vec{L}\perp\vec{p}=\mu\,\dot{\vec{r}}.
\]
En el plano orbital:
\[
\ell=\mu\,r^2\,\dot{\phi},\qquad
E=\frac{1}{2}\mu\big(\dot r^{\,2}+r^2\dot{\phi}^{\,2}\big)-\frac{\kappa}{r}.
\]
Si \(|\vec{r}\times\vec{p}|=0\) entonces \(r\,p\,\sin\angle(\hat{\mathbf r},\hat{\mathbf p})=0\), lo que implica velocidad puramente radial (caso \(L=0\)).

\subsection{Caso con plano $z=0$}
Sin pérdida de generalidad, en el plano orbital:
\[
\theta=\frac{\pi}{2},\qquad \dot{\theta}=0,\qquad \vec{L}=\ell\,\hat{\mathbf z}.
\]

\subsection{Lagrangiana útil a recordar (coordenadas polares en el plano)}
\[
L=\frac{1}{2}\mu\big(\dot r^{2}+r^{2}\dot{\phi}^{2}\big)+\frac{\kappa}{r},\qquad
\ell=\mu r^{2}\dot{\phi},\qquad
E=\frac{1}{2}\mu\big(\dot r^{2}+r^{2}\dot{\phi}^{2}\big)-\frac{\kappa}{r}.
\]

\subsection*{Órbitas cónicas y ecuación polar}
La solución general (leyes de Kepler) es una \textbf{cónica}:
\[
r(\phi)=\frac{p}{1+e\cos(\phi-\bar{\phi})},
\qquad
p=\frac{\ell^{2}}{\mu\,\kappa},\qquad
e=\sqrt{1+\frac{2E\ell^{2}}{\mu\,\kappa^{2}}}\ \ge 0.
\]
Equivalente en la forma que usas:
\[
\frac{1}{r}=\frac{\mu\kappa}{\ell^{2}}
\left[
1+\sqrt{1+\frac{2E\ell^{2}}{\mu\kappa^{2}}}\,\cos(\phi-\bar{\phi})
\right].
\]
Clasificación por energía: \(E<0\) (elipse), \(E=0\) (parábola), \(E>0\) (hipérbola).

\medskip
\emph{En cualquier área, si se hace con calidad no es cualquier cosa: la excelencia exige método y detalle.}

\subsection*{Fechas y logística (curso)}
Entrega de tareas: \textbf{martes 21}. \\
Posible examen: \textbf{lunes 28 de octubre}. \\
A partir de entonces: \textbf{ecuaciones de Hamilton}.

\subsection{Caso para determinar la geometría y el marco de referencia}

Para fijar la \textbf{geometría} (plano, ángulos y rotaciones) necesitamos:
\begin{itemize}
  \item \textbf{Condiciones iniciales} respecto a un marco de referencia inercial.
  \item El \textbf{momento angular orbital} \(\vec{L}\), que fija el plano de la órbita.
\end{itemize}

En el caso planar estándar, tomamos el \textbf{plano orbital} como \(z=0\), de modo que \(\vec{r}\in\mathbb{R}^2\) y \(\vec{L}=\ell\,\hat{\mathbf{z}}\).

\medskip

\noindent\textbf{Ángulo de orientación de la cónica.}  
La orientación de la cónica en el plano se describe por un ángulo de referencia:
\[
\phi \quad \text{o bien} \quad \bar{\phi},
\]
donde \(\bar{\phi}\) suele representar un \emph{desfase angular} (por ejemplo, ángulo del periastro respecto al eje \(x\) en la figura). En estas notas, \(\bar{\phi}\) es el ángulo que barre el eje mayor respecto al eje \(x\).

\begin{figure}[H]
    \centering
	    \includegraphics[width=0.7\textwidth]{5.png} % ajusta el ancho según convenga
    \caption{Descripción clara y concisa de la figura.}
    \label{fig:mi_figura2}
\end{figure}

\begin{figure}[H]
    \centering
	    \includegraphics[width=0.7\textwidth]{3.png} % ajusta el ancho según convenga
    \caption{Descripción clara y concisa de la figura.}
    \label{fig:mi_figura2}
\end{figure}

\noindent\textbf{Ecuación polar de la cónica en el plano \(z=0\).}  
Conservando tu notación (\(m\) masa reducida, \(\ell\) momento angular escalar, \(E\) energía), la órbita se escribe como
\[
\frac{1}{r} \;=\; \frac{m\,k}{\ell^{2}} \left[
  1 + \sqrt{\,1+\frac{2E\,\ell^{2}}{m\,k^{2}}\,}\,
  \cos\!\big(\phi - \bar{\phi}\big)
\right],
\]
donde el término bajo la raíz determina la \textbf{excentricidad} (\(E<0\Rightarrow\) elipse, \(E=0\Rightarrow\) parábola, \(E>0\Rightarrow\) hipérbola).  
\emph{Si la energía \(E\) cambia, cambia el tipo de trayectoria.}

\medskip

\noindent\textbf{Vector de posición en el plano orbital.}  
En el plano \(z=0\), el vector de posición es
\[
\vec{r}_{z=0} \;=\; r\big(\cos\phi\,\hat{\mathbf{i}}+\sin\phi\,\hat{\mathbf{j}}\big).
\]

\noindent\textbf{Rotación del sistema (cambio de base/orientación).}  
Si la cónica necesita rotarse un ángulo \(\alpha\) alrededor de \(\hat{\mathbf{z}}\) (por ejemplo, para alinear el periastro con el eje \(x\) o para introducir \(\bar{\phi}\)), usamos la matriz de rotación en \(2\)D o en \(3\)D:
\[
R_z(\alpha) \;=\;
\begin{pmatrix}
\cos\alpha & -\sin\alpha & 0\\
\sin\alpha & \phantom{-}\cos\alpha & 0\\
0 & 0 & 1
\end{pmatrix},
\qquad
\vec{r}_{np} \;=\; R_z(\alpha)\,\vec{r}.
\]
En tu notación compacta:
\[
\vec{r}_{np} \;=\; \big(R\,\vec{r}\big),
\]
donde \(R\) es la rotación adecuada (p. ej., \(R_z(\bar{\phi})\) para desplazar la cónica por \(\bar{\phi}\)).  
Si más adelante el plano orbital no coincidiera con \(z=0\), puede emplearse una \emph{secuencia de Euler} \((\Omega,i,\omega)\), pero para este curso basta con la rotación en el plano.

\medskip

\noindent\textbf{Descomposición CM/relativo (para ubicar cada cuerpo).}
\[
\vec{R} \;=\; \frac{ m_{1}\,\vec{r}_{1} + m_{2}\,\vec{r}_{2} }{ m_{1}+m_{2} },
\qquad
\vec{r} \;=\; \vec{r}_{2}-\vec{r}_{1}.
\]
A partir de esto, las posiciones de cada cuerpo en función de \(\vec{R}\) y \(\vec{r}\) son

\[  \vec{r}_{1} = \vec{r}_{1}(\vec{R}, \vec{r}) \]

\[ \vec{r}_{2} = \vec{r}_{2} (\vec{R}, \vec{r})  \]


\[
\boxed{
\ \vec{r}_{1} \;=\; \vec{R} \;-\; \frac{m_{2}}{m_{1}+m_{2}}\,\vec{r},
\qquad
\vec{r}_{2} \;=\; \vec{R} \;+\; \frac{m_{1}}{m_{1}+m_{2}}\,\vec{r}\ }.
\]
Estas relaciones son las que luego se \emph{rotan} con \(R\) si deseas expresar la órbita en un marco girado.

\begin{figure}[H]
    \centering
	    \includegraphics[width=0.7\textwidth]{4.png} % ajusta el ancho según convenga
    \caption{Descripción clara y concisa de la figura.}
    \label{fig:mi_figura2}
\end{figure}

\noindent\textbf{Resumen operativo para elegir geometría:}
\begin{enumerate}
  \item Fija el \textbf{plano orbital} con \(\vec{L}\) (tomar \(z=0\) es natural si \(\vec{L}\parallel\hat{\mathbf{z}}\)).
  \item Usa la ecuación polar \(r(\phi)\) y el \textbf{desfase} \(\bar{\phi}\) para orientar la cónica.
  \item Construye \(\vec{r}_{z=0}\) y, si hace falta, aplica \(\vec{r}_{np}=R_z(\alpha)\vec{r}\).
  \item Recupera \(\vec{r}_{1},\vec{r}_{2}\) con \(\vec{R}\) y \(\vec{r}\) (y rota si el marco lo exige).
\end{enumerate}

\medskip

\noindent\textbf{Nota práctica (código/Python).}  
Para graficar: (i) genera un arreglo de \(\phi\), (ii) evalúa \(r(\phi)\), (iii) crea \(\vec{r}_{z=0}\), (iv) aplica rotación \(R_z(\bar{\phi})\) si deseas alinear con tu figura, y (v) si necesitas la posición de cada cuerpo, usa las fórmulas de \(\vec{r}_1,\vec{r}_2\).  
En caso de dudas, consultar al profesor en su cubículo.

\subsection{Iniciamos el curso con la segunda ley de Newton}


\[ m \frac{  d^{2}\vec{r} }{dt^{2}} = \vec{F}_{a} + \vec{F}_{c} \]

imaginemos una partícula libre de moberse con 3 grados de libertad en el espacio $3D$.


Ahora imaginemos que la partícula sigue en el mismo plano $3D$, pero ahora confinada dentro de un cilíndro.


Supongamos que la particula puede moverse alrededor de la superficie del cilindro si fuese libre, tendríamos la siguiete ecuación:

Osea no todos los puntos son accesibles de la particula, solo los de la superficie de la esfera

\[ \vec{r} = x\hat{x} + y\hat{y} + z\hat{z}   \]

\[ \boxed{x^{2} + y^{2} = a^{2}}   \] 

\[ x=\rho \cos \phi, \quad y=\rho\sen\phi  \]

\[ \vec{r} = a \left( \cos\phi \hat{x} + \sen\phi\hat{j}   \right) + z\hat{z}   \]


\[  \boxed{ \left( m \frac{d^{2} \vec{r}}{dt^{2}}  \right) \cdot \delta{\vec{r}} = \vec{F}_{c} \cdot \vec{r}    } \]



Dalamberte dijo que las fuerzas de constriccion no te las dan, se pueden obtener de difernte manera. Dice, agarra M y derivada segunda de r depescto al tiempo, si es libre eso es igual a la fuerza de constricción que no conocemos. Proyecta eso a lo largo de los deplaxamineot que son lrmitidos para la particula. Es decir:





\[ \delta{\vec{r}}  = \frac{ \partial \vec{r}}{\partial \phi} \delta{\phi} + \frac{ \partial \vec{r}  }{\partial z} \delta{z}\]

La fuerza de constricción la conocemos hasta que resolvemos el problema.


\[ \boxed{m \frac{d^{2} \vec{r} }{ dt^{2}  } = \vec{F}_{c}  } \quad a \longleftrightarrow \rho  \]

Eso es cierto solo sobre esa ecuaciónd e consstricción, sobre $\rho=a$

¿Qué hacemos en cambiar de $a$ a $\rho$? es un problema en esencia distinto.

\end{document}